\documentclass[10pt,twocolumn]{article}
\usepackage[T1]{fontenc}
\usepackage[utf8]{inputenc}
\usepackage[a4paper, top=1.8cm, bottom=1.8cm, left=1.8cm, right=1.8cm]{geometry}
\usepackage{times}
\usepackage{booktabs}
\usepackage{array}
\usepackage{multirow}
\usepackage{graphicx}
\usepackage{amsmath}
\usepackage{hyperref}
\usepackage{url}
\usepackage{microtype}
\usepackage{xcolor}
\usepackage{caption}
\usepackage{float}
\usepackage{balance}
\usepackage{tabularx}
\usepackage{enumitem}

\captionsetup{font=small, labelfont=bf}
\setlength{\columnsep}{0.6cm}
\hypersetup{colorlinks=true, linkcolor=black, citecolor=black, urlcolor=black}

\title{
  \large \textbf{LinguoMT: A Zero-Shot Benchmark for African Language\\[0.2em]
  Speech Translation and Recognition}\\[0.5em]
  \normalsize \textit{Anonymous Submission --- Double-Blind Review}\\[0.3em]
  \small \textit{[DEBUG mode results --- indicative only; full-scale evaluation in preparation]}
}
\author{}
\date{}

\begin{document}
\maketitle
\thispagestyle{empty}

\begin{abstract}
We present \textbf{LinguoMT}, a systematic zero-shot evaluation of
state-of-the-art multilingual speech models on four African languages
spanning two architectures and two independent datasets.
We evaluate \texttt{facebook/seamless-m4t-v2-large} (end-to-end)
and a Whisper-large-v3 + NLLB-200 cascade (pipeline) on Yoruba,
Hausa, Igbo, and Swahili across the FLEURS benchmark and the
African-Celtic dataset.
In DEBUG mode (8 text samples, 3 audio samples per language),
SeamlessM4T-v2 achieves BLEU \textbf{30.70} for Igbo$\rightarrow$English
and \textbf{14.89} for Yoruba$\rightarrow$English on gold transcripts
from the African-Celtic dataset, confirming strong decoder capability
on clean input.
ASR word error rates reach \textbf{101.4\%} for Igbo and
\textbf{101.1\%} for Yoruba under SeamlessM4T-v2, indicating
hallucination rather than transcription failure, while Swahili ---
a non-tonal Bantu language --- achieves WER \textbf{47.4\%} under
both architectures, directly isolating lexical tone as the primary
acoustic encoding challenge.
On FLEURS, text translation BLEU collapses below~1.0 for all
African languages under both architectures, revealing a severe
domain gap between the controlled African-Celtic recordings and
FLEURS benchmark conditions.
These debug-mode findings establish directional baselines and
motivate the targeted adaptation, preprocessing, and transfer
strategies explored in companion papers of the LinguoMT series.
\end{abstract}

\textbf{Keywords:} African languages, speech translation, ASR,
FLEURS, SeamlessM4T, Whisper, NLLB-200, zero-shot evaluation,
low-resource, tonal languages, benchmark

\section{Introduction}

Speech is the primary communication mode for hundreds of millions
of African language speakers, yet access to reliable automatic
speech translation remains a privilege concentrated in a small
number of high-resource languages.
Three motivations drive this work.
First, African languages are systematically excluded from the
training pipelines of the very models that claim multilingual
coverage: published evaluations routinely report macro-averages
that conceal near-complete failure for individual low-resource
languages.
Second, the infrastructure cost of this exclusion is not merely
academic --- it affects healthcare communication, legal access,
and educational technology for speaker communities numbering in
the tens of millions.
Third, the rapid pace of multilingual model development, from
SeamlessM4T~\cite{barrault2023seamless} to
Whisper~\cite{radford2022whisper} and NLLB-200~\cite{nllb2022},
has outpaced rigorous, reproducible evaluation on African languages,
leaving practitioners without reliable baselines for deployment decisions.

These circumstances give rise to three interrelated problems.
The first is the absence of controlled, comparable evaluation:
existing per-language results are scattered across papers with
heterogeneous preprocessing, metrics, and dataset splits,
making cross-paper comparison unreliable.
The second is the architecture selection problem: end-to-end
models such as SeamlessM4T-v2 and cascaded pipelines such as
Whisper+NLLB-200 make fundamentally different trade-offs,
but no study has compared them head-to-head on the same
African language data under identical conditions.
The third is the bottleneck identification problem: when a
speech translation system fails on a low-resource African
language, it is rarely clear whether the failure originates in
the acoustic encoder, the translation decoder, or the interaction
between the two --- a distinction that carries direct implications
for where research investment should be directed.

From these problems, two research questions organise this study.
\textbf{RQ1:} How do end-to-end and cascaded architectures compare
in zero-shot performance across Yoruba, Igbo, Hausa, and Swahili,
and where does the primary performance bottleneck lie?
\textbf{RQ2:} What is the magnitude of the performance gap between
African languages and high-resource reference languages, and which
aspects are attributable to acoustic modelling versus translation quality?

To address these questions, this paper presents LinguoMT, a
systematic zero-shot benchmark evaluated across two datasets
and two architectures.
The experimental design isolates the acoustic encoder's contribution
by comparing audio-path and text-path performance under identical
decoders, and documents architecture-level coverage gaps that
preclude any single open system from serving all four target languages.
The remainder of the paper is structured as follows.
Section~2 surveys related work. Section~3 describes the system
architectures. Section~4 details the experimental setup.
Section~5 presents results. Section~6 discusses implications.
Section~7 concludes.

\section{Related Work}

\subsection{Multilingual Speech Foundation Models}

SeamlessM4T~\cite{barrault2023seamless} introduced a unified
architecture for speech and text translation across more than
100 languages, combining a W2v-BERT~2.0 acoustic encoder with
a multilingual sequence-to-sequence decoder trained on over
470,000 hours of aligned speech-text data.
Whisper~\cite{radford2022whisper} adopted a complementary
approach, training a standard encoder-decoder transformer on
680,000 hours of weakly supervised multilingual audio,
achieving competitive ASR across 99 languages without
task-specific fine-tuning.
NLLB-200~\cite{nllb2022} addressed the machine translation side,
training on 200 languages with an explicit equity mandate.
Pratap et al.~\cite{pratap2023mms} demonstrated that MMS-300M,
trained on over 1,100 languages, achieves competitive WER on
FLEURS for some African languages, though zero-shot performance
on low-resource tonal languages remains a persistent challenge.

Despite these advances, no study has evaluated these systems
under a controlled, reproducible protocol that compares
architectures head-to-head on the same African language data
with the same metrics.
\textbf{Remaining challenge~1:} A systematic, comparable
zero-shot benchmark for African languages that isolates
architectural contributions from data preparation artefacts
is absent from the literature.

\subsection{African Language Speech Resources}

Olatunji et al.~\cite{olatunji2023afrispeech} released
AfriSpeech-200, a pan-African clinical and general ASR
benchmark spanning 120 accents, demonstrating that accent
variation constitutes a fundamental bottleneck for models
deployed across African contexts.
Babu et al.~\cite{babu2022xlsr} showed that XLS-R achieves
competitive WER after fine-tuning on small African language
corpora, but zero-shot performance degrades substantially,
particularly for tonal languages with limited pre-training
representation.
Emezue et al.~\cite{emezue2025naijavoices} released NaijaVoices,
a large-scale corpus of culturally rich Nigerian speech, while
\`{O}g\'{u}nr\`{e}m\`{i} et
al.~\cite{ogunremi2024iroyinspeech} introduced
\`{I}roy\`{i}nSpeech to capture the tonal and diacritical
properties of Yoruba.
Maltais et al.~\cite{maltais2026naijas2st} provide the most
comprehensive recent cross-system evaluation on Nigerian
languages, reporting that zero-shot SeamlessM4T achieves
SSA-COMET scores around 20.6 for Igbo --- described as
near-random performance that underscores the magnitude of
the coverage gap.

\textbf{Remaining challenge~2:} No existing benchmark directly
quantifies the gap between African tonal languages and
high-resource references under identical zero-shot conditions.

\subsection{Cascade versus End-to-End Architectures}

Etchegoyhen et al.~\cite{etchegoyhen2022cascade} provide a
systematic case study concluding that cascaded systems remain
competitive when the ASR component is strong, but are
structurally vulnerable to error propagation.
End-to-end systems such as SeamlessM4T bypass this bottleneck
in principle, but the acoustic encoder must bear the full burden
of acoustic modelling for every supported language.

\textbf{Remaining challenge~3:} For low-resource African languages,
the relative contribution of the ASR bottleneck versus translation
quality to overall failure has not been systematically isolated
by comparing audio-path and text-path performance under an
identical decoder on the same data.
This paper fills that gap directly.

\section{System Overview}

\subsection{SeamlessM4T-v2 Large}

We evaluate \texttt{facebook/seamless-m4t-v2-large}, a
2.3-billion-parameter end-to-end model trained on the
SeamlessAlign dataset spanning over 470,000 hours of aligned
speech-text pairs.
For speech-to-text translation, the model receives raw audio
in a single forward pass through the W2v-BERT~2.0 acoustic
encoder, followed by a multilingual decoder that generates
English text directly.
For ASR evaluation, source and target language codes are set
identically.
For text MT evaluation, the audio encoder is bypassed entirely
and the model operates on gold transcripts, isolating the
translation decoder's performance from acoustic processing.
Speech-supported languages are Yoruba (\texttt{yor}),
Igbo (\texttt{ibo}), and Swahili (\texttt{swh}).
Hausa (\texttt{hau}) supports text-only translation but is
absent from the speech output vocabulary and is excluded from
S2TT and ASR evaluations.

\subsection{Whisper-large-v3 + NLLB-200 Cascade}

The cascade pipeline combines \texttt{openai/whisper-large-v3}
for ASR with \texttt{facebook/nllb-200-distilled-600M} for
translation into English.
Igbo is excluded from the cascade because Whisper has no
dedicated Igbo language token.
Supported languages are Yoruba (\texttt{yo}/\texttt{yor\_Latn}),
Hausa (\texttt{ha}/\texttt{hau\_Latn}), and Swahili
(\texttt{sw}/\texttt{swh\_Latn}).

\subsection{Architecture Comparison}

Table~\ref{tab:arch} summarises the key properties of the two
systems.
The complementary coverage gap is itself a substantive finding:
no single open architecture currently provides speech translation
coverage for all four target languages simultaneously.

\begin{table}[H]
\centering
\caption{Architecture comparison.}
\label{tab:arch}
\small
\setlength{\tabcolsep}{3pt}
\begin{tabular}{lll}
\toprule
\textbf{Property} & \textbf{SeamlessM4T-v2} & \textbf{Whisper+NLLB} \\
\midrule
Parameters     & $\sim$2.3B        & $\sim$2.1B total \\
Architecture   & Unified E2E       & Sequential cascade \\
Igbo speech    & \checkmark        & $\times$ \\
Hausa speech   & $\times$          & \checkmark \\
GPU required   & $\geq$16GB VRAM   & $\geq$16GB VRAM \\
\bottomrule
\end{tabular}
\end{table}

\section{Experimental Setup}

\subsection{Datasets}

The \textbf{FLEURS} benchmark~\cite{conneau2022fleurs} provides
multilingual spoken audio aligned with text transcripts across
102 languages, recorded by native speakers at 16~kHz.
We use the validation split with language configurations
\texttt{yo\_ng} (Yoruba), \texttt{ha\_ng} (Hausa),
\texttt{ig\_ng} (Igbo), and \texttt{sw\_ke} (Swahili).

The \textbf{African-Celtic dataset}~\cite{mcgill2024african}
(\texttt{McGill-NLP/african\_celtic\_dataset}) provides an
independent evaluation setting at 48~kHz sampling rate.
We use the dev split.
Hausa audio is absent from this dataset; Igbo is present but
excluded from the cascade pipeline due to the absence of a
Whisper Igbo language token.

All experiments reported here use \textbf{DEBUG mode}:
8 text samples and 3 audio samples per language.
This provides directional findings sufficient for benchmarking
architecture comparisons while remaining computationally
tractable.
Full-scale evaluation (300 samples per language) is in
preparation and reported in the companion extended version.

\subsection{Languages and Acoustic Properties}

Table~\ref{tab:languages} summarises the four target languages.
Table~\ref{tab:eda} reports the exploratory data analysis of the
African-Celtic audio samples, revealing that Yoruba exhibits
substantially higher silence ratio (0.485) and lower RMS energy
(0.052) than Igbo (silence 0.326, RMS 0.141), consistent with the
different recording environments and speaker prosodic patterns.

\begin{table}[H]
\centering
\caption{Target languages and dataset coverage.}
\label{tab:languages}
\small
\setlength{\tabcolsep}{3pt}
\begin{tabular}{lllcc}
\toprule
\textbf{Language} & \textbf{Family} & \textbf{Subfamily}
  & \textbf{FLEURS} & \textbf{AC} \\
\midrule
Yoruba  & Niger-Congo  & Volta-Niger & \checkmark & \checkmark \\
Hausa   & Afro-Asiatic & Chadic      & \checkmark & $\times$ \\
Igbo    & Niger-Congo  & Volta-Niger & \checkmark & \checkmark \\
Swahili & Niger-Congo  & Bantu       & \checkmark & $\times$ \\
\bottomrule
\multicolumn{5}{l}{\footnotesize AC = African-Celtic dataset.}
\end{tabular}
\end{table}

\begin{table}[H]
\centering
\caption{EDA: acoustic properties of African-Celtic dev audio.}
\label{tab:eda}
\small
\setlength{\tabcolsep}{3pt}
\begin{tabular}{lrrr}
\toprule
\textbf{Language} & \textbf{Duration (s)} & \textbf{Silence ratio} & \textbf{RMS} \\
\midrule
Igbo    & 8.70  & 0.326 & 0.141 \\
Yoruba  & 12.55 & 0.485 & 0.052 \\
\midrule
Mean    & 10.60 & 0.405 & 0.097 \\
\bottomrule
\end{tabular}
\end{table}

\subsection{Metrics}

All metrics are computed with \texttt{sacrebleu}~\cite{post2018sacrebleu}
and \texttt{jiwer}.
BLEU is computed with tokenisation \texttt{13a}; chrF uses
$\beta=2$.
Word error rate (WER) and character error rate (CER) are
computed for ASR evaluation.
WER above 1.0 indicates net word insertion (hallucination).

\section{Results}

\subsection{Text Translation --- African-Celtic}
\label{sec:ac_text}

Table~\ref{tab:ac_text} reports text translation on the
African-Celtic dev split using gold transcripts as input,
isolating the translation decoder from acoustic processing.

\begin{table}[H]
\centering
\caption{Text translation BLEU / chrF --- African-Celtic
(8 samples, DEBUG mode). Gold transcripts. Lower number precision reflects small sample size.}
\label{tab:ac_text}
\small
\setlength{\tabcolsep}{3pt}
\begin{tabular}{llrr}
\toprule
\textbf{Model} & \textbf{Lang.} & \textbf{BLEU} & \textbf{chrF} \\
\midrule
\multicolumn{4}{l}{\textit{Source $\rightarrow$ English}} \\
SeamlessM4T-v2 & Igbo   & \textbf{30.70} & 56.35 \\
SeamlessM4T-v2 & Yoruba & 14.89          & 42.46 \\
Whisper+NLLB   & Yoruba & 12.72          & 43.44 \\
Whisper+NLLB   & Hausa  & \textbf{30.31} & 54.67 \\
\midrule
\multicolumn{4}{l}{\textit{English $\rightarrow$ Source}} \\
SeamlessM4T-v2 & Igbo   & 23.74          & 51.24 \\
SeamlessM4T-v2 & Yoruba & 2.72           & 22.12 \\
Whisper+NLLB   & Yoruba & 4.69           & 32.17 \\
Whisper+NLLB   & Hausa  & 16.17          & 46.71 \\
\bottomrule
\multicolumn{4}{l}{\footnotesize DEBUG: 8 samples. Results are directional.}
\end{tabular}
\end{table}

Igbo$\rightarrow$English achieves BLEU 30.70 under SeamlessM4T-v2,
confirming the decoder's capacity to handle Igbo when presented
with clean text.
Hausa$\rightarrow$English reaches BLEU 30.31 via Whisper+NLLB,
demonstrating comparable text MT quality across both architectures
when the acoustic bottleneck is removed.
The pronounced asymmetry in the reverse direction ---
English$\rightarrow$Yoruba reaches only BLEU 2.72 under
SeamlessM4T-v2 --- reflects the structural difficulty of generating
tonal diacritical output.
The relatively higher English$\rightarrow$Yoruba score under
Whisper+NLLB (4.69) compared to SeamlessM4T-v2 (2.72) may reflect
NLLB-200's more targeted African language training data, though the
small sample size warrants caution in interpreting this gap.

\subsection{Text Translation --- FLEURS}

Table~\ref{tab:fleurs_text} reports text translation on the FLEURS
validation split using gold transcripts as input.

\begin{table}[H]
\centering
\caption{Text translation BLEU / chrF --- FLEURS validation
(8 samples, DEBUG mode). Gold transcripts.}
\label{tab:fleurs_text}
\small
\setlength{\tabcolsep}{3pt}
\begin{tabular}{llrr}
\toprule
\textbf{Model} & \textbf{Lang.} & \textbf{BLEU} & \textbf{chrF} \\
\midrule
\multicolumn{4}{l}{\textit{Source $\rightarrow$ English}} \\
SeamlessM4T-v2 & Igbo    & 0.43 & 17.06 \\
SeamlessM4T-v2 & Yoruba  & 0.59 & 15.96 \\
SeamlessM4T-v2 & Swahili & 0.41 & 15.77 \\
Whisper+NLLB   & Yoruba  & 0.57 & 15.55 \\
Whisper+NLLB   & Hausa   & 0.46 & 17.43 \\
Whisper+NLLB   & Swahili & 0.45 & 15.85 \\
\midrule
\multicolumn{4}{l}{\textit{English $\rightarrow$ Source}} \\
SeamlessM4T-v2 & Igbo    & 0.43 & 13.40 \\
SeamlessM4T-v2 & Yoruba  & 0.31 &  9.51 \\
SeamlessM4T-v2 & Swahili & 0.45 & 14.92 \\
Whisper+NLLB   & Yoruba  & 0.20 & 10.73 \\
Whisper+NLLB   & Hausa   & 0.44 & 16.46 \\
Whisper+NLLB   & Swahili & 0.36 & 14.71 \\
\bottomrule
\multicolumn{4}{l}{\footnotesize DEBUG: 8 samples. Results are directional.}
\end{tabular}
\end{table}

The collapse from African-Celtic BLEU 30.70 to FLEURS BLEU 0.43
for Igbo is the most striking finding of the text translation
analysis.
Under identical models and gold transcript input, performance
drops by more than two orders of magnitude.
FLEURS utterances are derived from the FLoRes-200 text, which
contains longer and more syntactically complex sentences than
the controlled conversational recordings of the African-Celtic
dataset.
This domain gap directly addresses RQ2: a substantial portion
of the performance attributed to ``African language difficulty''
may in fact reflect benchmark sensitivity rather than fundamental
model limitation.

\subsection{ASR Performance}

Table~\ref{tab:asr} reports ASR on the African-Celtic dataset.
The Swahili results serve as a non-tonal control for RQ1.

\begin{table}[H]
\centering
\caption{ASR WER and CER --- African-Celtic (3 samples, DEBUG mode).
WER $>$ 1.0 indicates net word insertion. Lower is better.}
\label{tab:asr}
\small
\setlength{\tabcolsep}{3pt}
\begin{tabular}{llrr}
\toprule
\textbf{Model} & \textbf{Language} & \textbf{WER} & \textbf{CER} \\
\midrule
SeamlessM4T-v2 & Igbo    & 1.014 & 0.524 \\
SeamlessM4T-v2 & Yoruba  & 1.011 & 0.410 \\
SeamlessM4T-v2 & Swahili & \textbf{0.474} & \textbf{0.198} \\
\midrule
Whisper+NLLB   & Igbo    & 0.797 & 0.369 \\
Whisper+NLLB   & Yoruba  & 0.863 & 0.342 \\
Whisper+NLLB   & Swahili & \textbf{0.474} & \textbf{0.198} \\
\bottomrule
\multicolumn{4}{l}{\footnotesize DEBUG: 3 samples. Hausa: no speech encoder (WER=1.0, omitted).}
\end{tabular}
\end{table}

WER above 1.0 for Igbo and Yoruba under SeamlessM4T-v2 indicates
not merely transcription failure but active hallucination: the
model inserts words at a rate exceeding the reference length.
The contrast with Swahili WER 0.474 --- a non-tonal Bantu language
--- is the single most important finding of the ASR analysis.
Both models achieve identical Swahili WER (0.474), suggesting the
acoustic encoder difficulty is not a generic African language
problem but is specifically tied to lexical tone encoding.
This directly answers RQ1.

\subsection{The Direction of the Gap: A Quantitative Summary}

Table~\ref{tab:summary} consolidates the key performance gaps
identified across experiments.

\begin{table}[H]
\centering
\caption{Performance gap summary (DEBUG mode).}
\label{tab:summary}
\small
\setlength{\tabcolsep}{3pt}
\begin{tabular}{p{4cm}rr}
\toprule
\textbf{Contrast} & \textbf{Tonal} & \textbf{Non-tonal} \\
\midrule
Igbo text BLEU (AC)             & 30.70 & --- \\
Igbo text BLEU (FLEURS)         & 0.43  & --- \\
Domain gap ratio                & $\times$71 & --- \\
\midrule
SeamlessM4T WER (tonal/non-tonal)  & 1.01--1.01  & 0.47 \\
Whisper WER (tonal/non-tonal)      & 0.80--0.86  & 0.47 \\
\bottomrule
\end{tabular}
\end{table}

\section{Discussion}

\subsection{Lexical Tone as the Primary Acoustic Bottleneck}

The ASR results resolve RQ1 with precision.
The Swahili non-tonal control demonstrates that acoustic encoder
failure is not a property of African languages in general, but of
tonal African languages specifically.
SeamlessM4T-v2 achieves WER 0.474 for Swahili while producing WER
above 1.0 for both Igbo and Yoruba --- languages that share
Niger-Congo family membership with Swahili but carry lexical
meaning through pitch contours.
The W2v-BERT~2.0 encoder, trained predominantly on non-tonal
speech, lacks the representational capacity to reliably encode
the fundamental frequency trajectories that distinguish lexical
items in tonal languages.
Whisper's lower (though still high) WER for tonal languages may
reflect its larger and more diverse pre-training corpus, though
the gap from Swahili to tonal languages remains large under both
architectures.

For practitioners, the implication is unambiguous: the translation
quality problem for Igbo and Yoruba is an acoustic representation
problem that requires targeted encoder adaptation rather than MT
model development.
NaijaS2ST~\cite{maltais2026naijas2st} provides direct confirmation,
reporting that monolingual fine-tuning of SeamlessM4T raises its
SSA-COMET score for Igbo from near-random (20.6 zero-shot) to 52.4
after acoustic adaptation.

\subsection{The African-Celtic to FLEURS Domain Gap}

The collapse of text translation BLEU from 30.70 on African-Celtic
to 0.43 on FLEURS for Igbo --- under identical models and gold
transcripts --- reveals a domain gap whose magnitude has not
previously been documented for these languages.
FLEURS utterances, derived from the FLoRes-200 sentence corpus,
average approximately 17.5 words per sentence (as measured by our
evaluation), against roughly 20 words for the African-Celtic dev
split.
The difference is not length alone: FLoRes sentences are drawn
from Wikipedia and news text, producing complex subordinate
clauses and domain-specific vocabulary that challenges the NLLB
decoder trained on more general-purpose parallel data.

This finding directly addresses RQ2: a substantial portion of the
performance gap attributed to African language difficulty may in
fact reflect the interaction between model capability and benchmark
domain rather than fundamental language-specific limitations.
Future evaluations should systematically report results across
multiple benchmarks to distinguish model quality from benchmark
sensitivity.

\subsection{Architecture Coverage as a Structural Finding}

The complementary coverage gap documented in Section~3 is not an
incidental limitation: it is a structural finding about the current
state of multilingual speech technology.
SeamlessM4T-v2 supports Igbo but not Hausa in speech mode; the
Whisper+NLLB cascade supports Hausa but not Igbo.
No open, freely available architecture currently covers all four
target languages in speech mode.
This fragmentation of coverage, across a set of languages
collectively spoken by over 200 million people, represents a
systemic failure of the universal multilingual model paradigm
that the research community has not yet adequately acknowledged.
Maltais et al.~\cite{maltais2026naijas2st} document the
magnitude of the Hausa gap explicitly, noting that the only
available workaround is to route Hausa audio through an Arabic
language proxy --- a strategy that produces near-random outputs
and underscores the absence of any meaningful Hausa acoustic
representation in current foundation models.

\subsection{Limitations and Caveats}

The DEBUG mode sample sizes (8 text, 3 audio per language)
are sufficient for directional findings but introduce substantial
variance on absolute metric values.
Full-mode evaluation using up to 300 samples per language is
required to produce publication-quality confidence intervals.
The identical Swahili WER of 0.474 under both architectures,
in particular, should be treated as a directional finding rather
than a definitive equivalence claim at this sample size.
Future work should extend evaluation to include European reference
languages (French, German, Spanish) evaluated under identical
conditions to establish the high-resource performance ceiling,
and should adopt SSA-COMET~\cite{li2025ssacomet} as a
complementary metric calibrated specifically for African MT outputs.

\section{Conclusion}

This paper has presented LinguoMT, a reproducible zero-shot
benchmark for African language speech translation and recognition
spanning two architectures and two datasets.
The central finding directly addresses RQ1: lexical tone, not
African language membership in general, is the primary acoustic
encoding bottleneck.
The Swahili non-tonal control achieves WER 0.474 under both
architectures while tonal Igbo and Yoruba reach WER above 0.80,
demonstrating that the encoder failure is tonality-specific.
The benchmark additionally documents a previously unreported
domain gap between the African-Celtic and FLEURS benchmarks,
and a structural architecture coverage gap that leaves no single
open system serving all four target languages.
These debug-mode findings motivate the adaptation strategies,
preprocessing interventions, cascade analyses, and cross-lingual
transfer experiments explored in the companion papers of the
LinguoMT series.

\balance
\bibliographystyle{plain}
\begin{thebibliography}{99}

\bibitem{babu2022xlsr}
A.~Babu et al.,
``XLS-R: Self-supervised Cross-lingual Speech Representation
Learning at Scale,''
in \textit{Proc.\ Interspeech}, 2022.

\bibitem{barrault2023seamless}
L.~Barrault et al.,
``SeamlessM4T: Massively Multilingual \& Multimodal Machine
Translation,''
\textit{arXiv:2308.11596}, 2023.

\bibitem{conneau2022fleurs}
A.~Conneau et al.,
``FLEURS: Few-Shot Learning Evaluation of Universal
Representations of Speech,''
in \textit{Proc.\ IEEE SLT}, pp.~798--805, 2022.

\bibitem{emezue2025naijavoices}
C.~Emezue et al.,
``The NaijaVoices Dataset,''
in \textit{Proc.\ Interspeech}, pp.~1338--1342, 2025.

\bibitem{etchegoyhen2022cascade}
T.~Etchegoyhen et al.,
``Cascade or Direct Speech Translation? A Case Study,''
\textit{Applied Sciences}, vol.~12, no.~3, 2022.

\bibitem{li2025ssacomet}
S.~Li et al.,
``SSA-COMET: Do LLMs Outperform Learned Metrics in Evaluating MT
for Under-Resourced African Languages?''
in \textit{Proc.\ EMNLP}, pp.~12979--12998, 2025.

\bibitem{maltais2026naijas2st}
M.~Maltais et al.,
``NaijaS2ST: A Multi-Accent Benchmark for Speech-to-Speech
Translation in Low-Resource Nigerian Languages,''
\textit{arXiv:2604.16287}, 2026.

\bibitem{mcgill2024african}
McGill-NLP,
``African-Celtic Dataset,''
\textit{Hugging Face}, 2024.

\bibitem{nllb2022}
NLLB Team,
``No Language Left Behind: Scaling Human-Centered Machine
Translation,''
\textit{arXiv:2207.04672}, 2022.

\bibitem{ogunremi2024iroyinspeech}
T.~\`{O}g\'{u}nr\`{e}m\`{i} et al.,
``\`{I}roy\`{i}nSpeech: A Multi-Purpose Yor\`{u}b\'{a}
Speech Corpus,''
in \textit{Proc.\ LREC-COLING}, pp.~9296--9303, 2024.

\bibitem{olatunji2023afrispeech}
T.~Olatunji et al.,
``AfriSpeech-200: Pan-African Accented Speech Dataset,''
\textit{Trans.\ ACL}, vol.~11, pp.~1669--1685, 2023.

\bibitem{post2018sacrebleu}
M.~Post,
``A Call for Clarity in Reporting BLEU Scores,''
in \textit{Proc.\ WMT}, 2018.

\bibitem{pratap2023mms}
V.~Pratap et al.,
``Scaling Speech Technology to 1,000+ Languages,''
\textit{arXiv:2305.13516}, 2023.

\bibitem{radford2022whisper}
A.~Radford et al.,
``Robust Speech Recognition via Large-Scale Weak
Supervision,''
in \textit{Proc.\ ICML}, 2023.

\end{thebibliography}

\end{document}
